\chapter{좋은 VLA 연구 주제를 고르는 법}
\label{ch:research_topic_selection}

13장에서는 VLA 연구 지형을 model, data, system, evaluation의 네 축으로 정리했다. 이 장은 그 지도를 연구 주제 선택 기준으로 바꾼다. 좋은 VLA 주제는 최신 모델을 가져와 새로운 task에 돌려보는 것만으로 만들어지지 않는다. 좋은 주제는 명확한 병목을 잡고, 그 병목이 왜 중요한지 설명하며, 방법이 그 병목을 직접 겨냥하고, 실험이 그 주장을 공정하게 검증할 때 만들어진다.

\section{좋은 주제의 기본 조건}

VLA 연구 주제는 다음 네 조건을 만족해야 한다.

\begin{enumerate}
  \item 병목이 명확해야 한다.
  \item contribution이 병목과 직접 연결되어야 한다.
  \item 평가가 contribution을 분리해서 측정해야 한다.
  \item 실패와 한계를 숨기지 않고 분석할 수 있어야 한다.
\end{enumerate}

이를 간단한 score로 쓰면 다음과 같다.

\begin{equation}
  Q_{\mathrm{topic}}
  =
  w_b B
  +
  w_n N
  +
  w_v V
  +
  w_f F
  -
  w_c C.
\end{equation}

$B$는 bottleneck clarity, $N$은 novelty, $V$는 verifiability, $F$는 feasibility, $C$는 confound risk이다. 좋은 주제는 새로워 보이는 것보다 검증 가능한 것이 중요하다. 특히 VLA는 모델, 데이터, controller, benchmark가 얽혀 있어 confound가 크다.

\section{피해야 할 주제 유형}

가장 피해야 할 유형은 ``최신 VLA를 우리 task에 적용했다''이다. 이 자체는 engineering report일 수 있지만, 연구 contribution은 약하다. 또 다른 위험한 유형은 benchmark만 바꾸고 원인을 설명하지 않는 연구이다.

\begin{table}[t]
  \centering
  \caption{약한 VLA 주제와 보강 방향}
  \label{tab:weak_vla_topics}
  \begin{tabularx}{\textwidth}{>{\bfseries}p{34mm}X}
    \toprule
    약한 주제 & 보강 방향 \\
    \midrule
    최신 모델 적용 & 실패 mode를 분석하고 그 mode를 줄이는 방법을 제안한다 \\
    데이터만 추가 & 어떤 data mixture가 어떤 일반화 축을 개선하는지 분리한다 \\
    새 benchmark 제안 & baseline, split, failure taxonomy, reproducibility를 함께 제공한다 \\
    속도 개선 주장 & success-latency-memory trade-off curve를 제시한다 \\
    reasoning 주장 & reasoning trace가 아니라 recovery success와 safety 개선을 보인다 \\
    safety 주장 & safety envelope와 risk metric을 실제 episode에서 검증한다 \\
    \bottomrule
  \end{tabularx}
\end{table}

VLA 분야에서는 demo video가 강한 설득력을 갖지만, 논문 contribution을 대신하지는 못한다. Video는 qualitative evidence이고, 연구 주장은 controlled comparison으로 증명해야 한다.

\section{병목에서 시작하기}

좋은 연구 질문은 모델 이름이 아니라 병목에서 시작한다. 예를 들어 다음은 모두 좋은 출발점이다.

\begin{itemize}
  \item Action chunk가 길어질 때 접촉 단계에서 failure가 증가한다.
  \item Offline action loss는 낮아지지만 real robot success는 오르지 않는다.
  \item Language paraphrase에는 강하지만 spatial relation에는 약하다.
  \item Faster decoding은 평균 latency를 줄이지만 tail latency와 failure를 늘린다.
  \item VLA verifier가 unsafe instruction은 잡지만 scene-specific risk는 놓친다.
\end{itemize}

이런 병목은 측정 가능하다. 측정 가능한 병목은 method와 experiment를 자연스럽게 만든다.

\begin{equation}
  \mathrm{Claim}
  =
  \mathrm{Bottleneck}
  +
  \mathrm{Mechanism}
  +
  \mathrm{Evidence}.
\end{equation}

이 식은 논문 기획의 기본이다. 병목만 있으면 문제 제기이고, mechanism만 있으면 방법 소개이며, evidence가 없으면 주장이 아니다.

\section{Contribution type 선택}

VLA 논문의 contribution은 여러 유형이 있다.

\begin{table}[t]
  \centering
  \caption{VLA contribution 유형}
  \label{tab:vla_contribution_types}
  \begin{tabularx}{\textwidth}{>{\bfseries}p{30mm}X}
    \toprule
    유형 & 좋은 contribution의 조건 \\
    \midrule
    Architecture & 기존 구조의 어떤 병목을 줄이는지 ablation으로 보여야 한다 \\
    Action representation & reconstruction이 아니라 closed-loop success와 failure mode를 보여야 한다 \\
    Fine-tuning method & data efficiency, forgetting, task adaptation을 분리해 측정해야 한다 \\
    Dataset & 새로운 일반화 축과 evaluation protocol을 명확히 해야 한다 \\
    Benchmark & 공정한 baseline, hidden split, failure taxonomy를 포함해야 한다 \\
    System & latency, controller, safety, real robot constraint를 end-to-end로 보고해야 한다 \\
    Theory/analysis & 실제 VLA 실험 해석에 도움 되는 quantity를 제공해야 한다 \\
    \bottomrule
  \end{tabularx}
\end{table}

학생 프로젝트에서는 architecture contribution보다 evaluation-centric 또는 system-centric contribution이 현실적일 때가 많다. 큰 foundation model을 처음부터 학습하기 어렵기 때문이다. 반대로 강한 실험실 infrastructure가 있다면 real-robot failure analysis와 robust fine-tuning이 좋은 주제가 될 수 있다.

\section{Baseline을 먼저 설계하기}

VLA 연구에서 baseline은 나중에 붙이는 비교 대상이 아니다. Baseline 설계가 주제의 품질을 결정한다. 좋은 baseline은 다음 조건을 갖는다.

\begin{enumerate}
  \item 같은 data budget을 사용한다.
  \item 같은 observation과 controller를 사용한다.
  \item 같은 evaluation split을 사용한다.
  \item method가 주장하는 component만 다르다.
\end{enumerate}

Ablation은 다음처럼 contribution을 분해한다.

\begin{equation}
  \Delta J_{\mathrm{method}}
  =
  J(\pi_{\mathrm{ours}})
  -
  J(\pi_{\mathrm{baseline}}).
\end{equation}

하지만 이 차이가 method 때문인지 data 때문인지 controller 때문인지 모르면 논문은 약해진다. 따라서 실험 설계의 목적은 높은 숫자를 얻는 것이 아니라 $\Delta J$의 원인을 설명하는 것이다.

\section{실험 규모와 현실성}

좋은 주제는 feasible해야 한다. VLA는 대규모 GPU, robot hardware, data collection이 필요한 경우가 많다. 따라서 주제를 고를 때 resource budget을 명시해야 한다.

\begin{equation}
  R_{\mathrm{budget}}
  =
  (G, T_{\mathrm{robot}}, N_{\mathrm{episodes}}, H_{\mathrm{human}}, D_{\mathrm{storage}}).
\end{equation}

$G$는 GPU 자원, $T_{\mathrm{robot}}$은 robot 시간, $N_{\mathrm{episodes}}$는 실행 episode 수, $H_{\mathrm{human}}$은 annotation 또는 teleoperation 시간, $D_{\mathrm{storage}}$는 data 저장 비용이다. 좋은 연구 계획은 이 budget 안에서 검증 가능한 claim을 고른다.

실제 수업 프로젝트라면 다음 스케일이 현실적이다.

\begin{itemize}
  \item Offline VLA 분석: 공개 dataset과 pretrained model을 사용해 action representation 또는 failure taxonomy를 분석한다.
  \item Small real-robot study: 3--5개 task, 3--5개 object, 20--50 episode 수준에서 명확한 hypothesis를 검증한다.
  \item Benchmark/tooling project: 기존 논문의 evaluation protocol을 재현하고 OOD split 또는 failure annotation을 추가한다.
  \item Efficient inference project: latency와 success trade-off를 동일 hardware에서 측정한다.
\end{itemize}

\section{논문 제목으로 바꾸기}

좋은 주제는 제목으로 바꿨을 때 contribution이 드러난다.

\begin{table}[t]
  \centering
  \caption{약한 제목과 강한 제목}
  \label{tab:vla_title_examples}
  \begin{tabularx}{\textwidth}{>{\bfseries}p{42mm}X}
    \toprule
    약한 제목 & 강한 제목 방향 \\
    \midrule
    Applying VLA to Household Tasks & Failure-Aware Fine-Tuning for Long-Horizon Household Manipulation \\
    Faster VLA Inference & Success-Preserving Action Chunking under Real-Time Robot Control Constraints \\
    A New VLA Benchmark & An OOD and Failure-Taxonomy Benchmark for Language-Guided Manipulation \\
    Reasoning for Robots & Grounded Replanning with Uncertainty-Aware Verifiers for VLA Policies \\
    Multi-Robot VLA & Embodiment-Conditioned Action Decoding for Cross-Robot VLA Transfer \\
    \bottomrule
  \end{tabularx}
\end{table}

강한 제목은 task, 병목, 방법, 평가 축 중 최소 두 가지를 포함한다. 제목이 강하다고 논문이 강해지는 것은 아니지만, 좋은 제목으로 압축되지 않는 주제는 대개 contribution이 흐릿하다.

\section{13-question framework로 주제 점검}

3장에서 제시한 13-question framework는 논문 읽기뿐 아니라 주제 기획에도 쓸 수 있다. 아직 논문이 없는 상태에서는 각 항목을 질문으로 뒤집는다.

\begin{itemize}
  \item 배경: 왜 지금 이 병목이 중요한가.
  \item 문제: 구체적으로 무엇이 실패하는가.
  \item 기존 한계: 기존 VLA가 왜 이 실패를 해결하지 못하는가.
  \item 목표: 어떤 metric을 얼마나 개선할 것인가.
  \item 방법: 어떤 mechanism으로 병목을 줄일 것인가.
  \item 검증: 어떤 dataset, robot, split, baseline으로 확인할 것인가.
  \item 한계: 어떤 조건에서는 여전히 실패할 것인가.
\end{itemize}

이 질문에 답하지 못하면 아직 연구 주제가 아니라 아이디어에 가깝다.

\section{요약}

좋은 VLA 연구 주제는 최신성을 따라가는 데서 나오지 않고, 명확한 병목과 검증 가능한 mechanism에서 나온다. Model-centric, data-centric, system-centric, evaluation-centric 중 어느 축에 서 있는지 정하고, baseline과 ablation을 먼저 설계하며, resource budget 안에서 증명 가능한 claim을 골라야 한다.

다음 장에서는 이 기준을 대표 논문 읽기에 적용한다. 모든 논문을 깊게 다룰 수는 없지만, VLA 연구를 처음 시작하는 학생이 반드시 읽어야 할 논문군을 유형별로 정리하고, 각 논문을 어떤 질문으로 읽어야 하는지 제시한다.

\begin{sourcebox}{Key papers for topic selection}
\begin{itemize}
  \item OpenVLA-OFT (2025): Claim = fine-tuning bottleneck + optimized decoding/action representation mechanism + LIBERO/real robot evidence로 읽기 좋다.
  \item FAST (2025): Claim = action tokenization bottleneck + DCT compression mechanism + high-frequency dexterity evidence로 읽기 좋다.
  \item SmolVLA (2025): Claim = deployment cost bottleneck + compact/asynchronous system mechanism + affordable hardware evidence로 읽기 좋다.
  \item VLA-OS/WorldVLA (2025): Claim = planning/world-model bottleneck + structured intermediate representation + planning utility evidence로 평가해야 한다.
  \item GR00T N1/Gemini Robotics (2025/2026): Claim = humanoid or closed industrial VLA + dual-model/dual-system mechanism + 공개 evidence의 한계로 읽어야 한다.
\end{itemize}
\end{sourcebox}

\begin{assignmentbox}{Reading assignment [Advanced]}
자신의 연구 아이디어를 \(\mathrm{Claim}=\mathrm{Bottleneck}+\mathrm{Mechanism}+\mathrm{Evidence}\) 형식으로 한 문장에 쓰고, 그 claim을 반박할 수 있는 baseline 두 개와 ablation 두 개를 설계하라.
\end{assignmentbox}
